% MATH companion editorial bridge VI/27; explanatory prose only.
\paragraph{From an integral domain to an integral scheme.}
An affine scheme \(\operatorname{Spec}A\) is integral exactly when \(A\)
is an integral domain.  The scheme-theoretic version of this condition
has two parts: the underlying topological space is irreducible, and the
structure sheaf has no nonzero nilpotents (the scheme is reduced).
Irreducibility alone does not rule out nilpotents; reducedness alone
does not prevent several irreducible components.

\paragraph{Why all nonempty affine opens have the same function field.}
Let \(X\) be an integral scheme and \(\eta\) its unique generic point.
Every nonempty open set contains \(\eta\): otherwise its complement,
a proper closed subset, would contain a dense point.
If \(U=\operatorname{Spec}A\) is any nonempty affine open, then
\(A\) is a domain and the local ring at the generic point is
\[
 \mathcal O_{X,\eta}\cong A_{(0)}=\operatorname{Frac}(A).
\]
This local ring does not depend on the choice of \(U\).
It is therefore legitimate to define the function field intrinsically by
\(K(X)=\mathcal O_{X,\eta}\).  The passage to the fraction field is not
an arbitrary choice of coordinates: it records precisely the germs of
rational functions at the generic point.

\paragraph{A calculation on two different affine charts.}
Take \(X=\mathbb P^1_k\).
On the usual affine chart \(U_0\) use the coordinate \(t=x_1/x_0\),
so that \(U_0\simeq\operatorname{Spec}k[t]\).
On the other affine chart \(U_1\) use \(s=x_0/x_1\),
so that \(U_1\simeq\operatorname{Spec}k[s]\).
The intersection is defined by \(t\neq 0\), where \(s=1/t\).
Hence the two presentations have fraction fields
\[
 \operatorname{Frac}k[t]=k(t)
 \qquad\text{and}\qquad
 \operatorname{Frac}k[s]=k(s),\quad s=t^{-1},
\]
which become the same field after the transition function is applied.
This explicit gluing calculation is a model for the general statement.

\paragraph{Dominant maps and the direction of function fields.}
For integral schemes, a dominant morphism \(f:X\to Y\) sends the generic
point of \(X\) to the generic point of \(Y\), and therefore induces a
field homomorphism
\[
 K(Y)\hookrightarrow K(X).
\]
Notice the reversal of direction: geometric maps pull rational
functions back.  For example, the map
\(\mathbb A^1_k\to\mathbb A^1_k\) given by \(t\mapsto t^2\)
induces \(k(u)\hookrightarrow k(t)\) by \(u\mapsto t^2\).
When the source is not integral, there may be no single generic point
and no single intrinsic function field of the above kind.
